\section{Giới thiệu}
\subsection{Động cơ của đề tài}

Trong bối cảnh tốc độ đô thị hóa và gia tăng dân số diễn ra mạnh mẽ trong nhiều năm gần đây, Thành phố Hồ Chí Minh (TP.HCM) đang phải đối mặt với những áp lực ngày càng lớn đối với hệ thống giao thông đô thị. Sự gia tăng nhanh chóng của phương tiện cá nhân, đặc biệt là xe máy và ô tô, trong khi năng lực hạ tầng giao thông chưa theo kịp tốc độ phát triển đô thị, đã dẫn đến tình trạng ùn tắc giao thông diễn ra thường xuyên tại nhiều tuyến đường và khu vực trọng điểm. Thực trạng này không chỉ làm suy giảm hiệu quả vận hành của đô thị mà còn gây ra nhiều hệ lụy nghiêm trọng như ô nhiễm môi trường, gia tăng thời gian di chuyển, tiêu hao năng lượng, tai nạn giao thông và ảnh hưởng trực tiếp đến chất lượng sống của người dân.

Trước yêu cầu phát triển đô thị bền vững, TP.HCM đã định hướng xây dựng mô hình đô thị thông minh, trong đó giao thông thông minh (Intelligent Transportation System – ITS) được xem là một trong những lĩnh vực ưu tiên trọng điểm. Tuy nhiên, phần lớn các hệ thống giám sát và quản lý giao thông hiện nay mới chỉ tập trung vào chức năng quan sát và phản ứng theo thời gian thực, trong khi còn thiếu một nền tảng dữ liệu tập trung có khả năng tích hợp, lưu trữ và phân tích dữ liệu giao thông trên quy mô lớn và dài hạn. Điều này làm hạn chế khả năng khai thác dữ liệu phục vụ công tác thống kê, đánh giá xu hướng, tối ưu hóa vận hành và hỗ trợ ra quyết định chiến lược.

Hiện nay, dữ liệu giao thông đô thị được hình thành từ nhiều nguồn dị biệt như hệ thống camera giám sát (CCTV), thiết bị định vị GPS, dữ liệu cảm biến giao thông, nền tảng bản đồ số, cũng như dữ liệu lưu lượng và sự cố được cung cấp từ các nền tảng giao thông trực tuyến. Các nguồn dữ liệu này có đặc trưng phân tán, không đồng nhất về cấu trúc và thiếu tính liên kết không gian – thời gian, dẫn đến khó khăn trong việc đồng bộ, tổng hợp và khai thác phân tích chuyên sâu. Đặc biệt, đối với các bài toán giao thông thông minh hiện đại, dữ liệu không gian địa lý (geospatial data) đóng vai trò then chốt trong việc mô hình hóa mạng lưới giao thông, xác định các tuyến trọng điểm, phân tích mật độ lưu thông và đánh giá mức độ ùn tắc theo từng khu vực.

Trong bối cảnh đó, mô hình kho dữ liệu giao thông (Traffic Data Warehouse) kết hợp với công nghệ dữ liệu không gian (Spatial Data Warehouse) được xem là giải pháp nền tảng nhằm giải quyết bài toán phân mảnh dữ liệu. Kho dữ liệu cho phép tích hợp dữ liệu từ nhiều nguồn khác nhau, chuẩn hóa và tổ chức dữ liệu theo cấu trúc đa chiều, hỗ trợ hiệu quả cho các truy vấn phân tích phức tạp, thống kê lịch sử và khai thác dữ liệu quy mô lớn. Khác với các hệ quản trị cơ sở dữ liệu giao dịch truyền thống (OLTP), kho dữ liệu được tối ưu hóa cho các tác vụ phân tích (OLAP), đặc biệt phù hợp với các bài toán xử lý dữ liệu giao thông theo chiều không gian – thời gian và phục vụ công tác điều hành đô thị thông minh.

Bên cạnh đó, sự phát triển mạnh mẽ của trí tuệ nhân tạo (Artificial Intelligence – AI) và học máy (Machine Learning) đang thúc đẩy quá trình chuyển đổi từ mô hình quản lý giao thông mang tính phản ứng (reactive traffic management) sang mô hình quản lý chủ động và dự báo (proactive traffic management). Các mô hình AI có khả năng phân tích xu hướng lưu lượng, dự báo nguy cơ ùn tắc và hỗ trợ tối ưu hóa điều phối giao thông theo thời gian thực. Tuy nhiên, hiệu quả của các mô hình dự báo phụ thuộc rất lớn vào chất lượng dữ liệu đầu vào, đòi hỏi một hạ tầng dữ liệu đủ lớn, ổn định, được chuẩn hóa và cập nhật liên tục. Đây chính là vai trò cốt lõi của hệ thống kho dữ liệu giao thông tích hợp.

Xuất phát từ những yêu cầu thực tiễn nêu trên, nhóm nghiên cứu quyết định thực hiện đề tài: “Mô hình kho dữ liệu giao thông phục vụ công tác thống kê, phân tích và dự báo hỗ trợ điều hành giao thông thông minh”. Đề tài tập trung xây dựng một nền tảng dữ liệu giao thông tích hợp trên cơ sở kết hợp dữ liệu thời gian thực từ các nguồn giao thông trực tuyến với dữ liệu không gian địa lý, triển khai quy trình ETL tự động nhằm thu thập, làm sạch, biến đổi và lưu trữ dữ liệu trong kho dữ liệu không gian. Đồng thời, hệ thống hướng đến việc ứng dụng các mô hình trí tuệ nhân tạo trong phân tích và dự báo tình trạng giao thông, kết hợp với nền tảng trực quan hóa và dashboard hỗ trợ giám sát, từ đó góp phần xây dựng nền tảng hạ tầng dữ liệu phục vụ phát triển hệ sinh thái giao thông thông minh và đô thị thông minh tại TP.HCM trong tương lai.


\subsection{Mục tiêu}

Mục tiêu tổng quát của đề tài là nghiên cứu, thiết kế và triển khai một nền tảng kho dữ liệu giao thông thông minh có khả năng tích hợp dữ liệu đa nguồn, xử lý dữ liệu không gian – thời gian theo thời gian thực và hỗ trợ các tác vụ phân tích, dự báo phục vụ công tác quản lý và điều hành giao thông đô thị. Hệ thống được định hướng xây dựng theo mô hình kiến trúc dữ liệu hiện đại, bảo đảm khả năng mở rộng, tính ổn định và đáp ứng yêu cầu khai thác dữ liệu quy mô lớn trong bối cảnh phát triển đô thị thông minh.

Để hiện thực hóa mục tiêu tổng quát nêu trên, đề tài tập trung vào các mục tiêu cụ thể sau:

\begin{itemize}

    \item \textbf{Nghiên cứu và đề xuất kiến trúc kho dữ liệu giao thông thông minh}\\
    Xây dựng kiến trúc tổng thể cho hệ thống kho dữ liệu giao thông theo mô hình đa tầng, bao gồm tầng thu thập dữ liệu, tầng lưu trữ dữ liệu không gian, tầng phân tích trí tuệ nhân tạo và tầng trực quan hóa dữ liệu. Đề tài tập trung thiết kế mô hình dữ liệu đa chiều phù hợp với đặc thù dữ liệu giao thông đô thị tại Việt Nam, kết hợp giữa dữ liệu cấu trúc, bán cấu trúc và dữ liệu không gian địa lý (geospatial data). Các mô hình Star Schema hoặc Snowflake Schema được nghiên cứu và áp dụng nhằm tối ưu hóa hiệu năng truy vấn, hỗ trợ hiệu quả cho các tác vụ phân tích dữ liệu lớn và xử lý theo chiều không gian – thời gian.

    \item \textbf{Xây dựng quy trình xử lý dữ liệu ETL/ELT tự động}\\
    Thiết kế và triển khai quy trình ETL/ELT nhằm tự động hóa quá trình thu thập, biến đổi và lưu trữ dữ liệu giao thông từ nhiều nguồn khác nhau như API giao thông thời gian thực, dữ liệu bản đồ số OpenStreetMap (OSM), dữ liệu sự cố giao thông và các nguồn dữ liệu hỗ trợ khác. Quy trình xử lý dữ liệu tập trung vào việc làm sạch dữ liệu, xử lý ngoại lệ, chuẩn hóa định dạng, tích hợp dữ liệu không gian và tối ưu khả năng cập nhật dữ liệu liên tục theo chu kỳ thời gian thực, bảo đảm tính toàn vẹn và độ tin cậy của dữ liệu trong kho dữ liệu trung tâm.

    \item \textbf{Xây dựng kho dữ liệu không gian phục vụ phân tích giao thông}\\
    Triển khai hệ quản trị cơ sở dữ liệu PostgreSQL kết hợp với phần mở rộng PostGIS nhằm hỗ trợ lưu trữ và xử lý dữ liệu không gian địa lý. Đề tài tập trung xây dựng mô hình dữ liệu phục vụ quản lý các thực thể giao thông như tuyến đường, đoạn đường, nút giao, vùng ảnh hưởng và dữ liệu lưu lượng theo không gian – thời gian. Đồng thời, hệ thống được tối ưu hóa nhằm đáp ứng các truy vấn phân tích không gian phức tạp và hỗ trợ trực quan hóa dữ liệu giao thông trên nền bản đồ số.

    \item \textbf{Tích hợp và triển khai mô hình trí tuệ nhân tạo phục vụ dự báo giao thông}\\
    Nghiên cứu và ứng dụng các mô hình học máy và học sâu trong bài toán dự báo tình trạng giao thông đô thị, bao gồm dự báo lưu lượng phương tiện, tốc độ lưu thông và nguy cơ ùn tắc. Các mô hình như ARIMA, LSTM hoặc các mô hình lai ghép (Hybrid Models) được huấn luyện và đánh giá trên tập dữ liệu lịch sử đã được chuẩn hóa từ kho dữ liệu giao thông. Quá trình huấn luyện và kiểm thử mô hình được thực hiện theo quy trình khoa học nhằm bảo đảm độ chính xác và khả năng ứng dụng thực tiễn của hệ thống dự báo.

    \item \textbf{Phát triển hệ thống trực quan hóa và hỗ trợ ra quyết định}\\
    Xây dựng nền tảng dashboard và web application phục vụ giám sát giao thông theo thời gian thực, hỗ trợ trực quan hóa dữ liệu dưới dạng biểu đồ, bản đồ nhiệt (heatmap), thống kê lưu lượng và cảnh báo ùn tắc. Hệ thống hướng đến việc cung cấp công cụ hỗ trợ ra quyết định cho các cơ quan quản lý giao thông đô thị, hỗ trợ công tác điều phối giao thông, phân luồng phương tiện và hoạch định chính sách phát triển hạ tầng giao thông thông minh.

    \item \textbf{Triển khai thử nghiệm và đánh giá hiệu năng hệ thống}\\
    Tiến hành triển khai thử nghiệm hệ thống trên hạ tầng điện toán đám mây Microsoft Azure, kết hợp với nền tảng container hóa Docker nhằm bảo đảm tính linh hoạt và khả năng mở rộng của hệ thống. Đề tài thực hiện đánh giá hiệu năng của quy trình xử lý dữ liệu thời gian thực, khả năng đáp ứng của kho dữ liệu, mức độ ổn định của hệ thống cũng như độ chính xác của mô hình dự báo trên dữ liệu thực tế. Các kết quả đánh giá được sử dụng làm cơ sở để phân tích tính khả thi và khả năng ứng dụng của hệ thống trong bài toán giao thông thông minh tại TP.HCM.

\end{itemize}

\subsection{Phạm vi}

Nhằm bảo đảm tính khả thi của đề tài trong điều kiện thời gian, nguồn lực và phạm vi nghiên cứu, hệ thống được giới hạn trong các nội dung nghiên cứu và triển khai cụ thể như sau:

\begin{itemize}

    \item \textbf{Về dữ liệu nghiên cứu}\\
    Đề tài tập trung khai thác và xử lý các loại dữ liệu giao thông có cấu trúc và bán cấu trúc phục vụ cho bài toán phân tích giao thông thông minh. Các nguồn dữ liệu chính bao gồm dữ liệu lưu lượng giao thông thời gian thực, dữ liệu tốc độ lưu thông, dữ liệu sự cố giao thông, dữ liệu định vị GPS và dữ liệu bản đồ số từ nền tảng OpenStreetMap (OSM). Ngoài ra, hệ thống có khả năng mở rộng để tích hợp dữ liệu từ các cảm biến IoT hoặc nền tảng giám sát giao thông trong tương lai. Đối với dữ liệu phi cấu trúc như video thô từ camera CCTV, đề tài không thực hiện xử lý trực tiếp mà chỉ sử dụng các dữ liệu đã được trích xuất và chuyển đổi sang dạng số liệu phục vụ phân tích.

    \item \textbf{Về phạm vi không gian nghiên cứu}\\
    Hệ thống được triển khai thử nghiệm và đánh giá trên một số tuyến đường trọng điểm, các hành lang giao thông chính và các khu vực có mật độ lưu thông cao tại Thành phố Hồ Chí Minh. Phạm vi nghiên cứu tập trung vào các tuyến đường thường xuyên xảy ra ùn tắc giao thông nhằm phục vụ việc phân tích lưu lượng, đánh giá mức độ quá tải và thử nghiệm khả năng dự báo tình trạng giao thông theo thời gian thực. Kết quả thử nghiệm được sử dụng làm cơ sở đánh giá tính khả thi và khả năng mở rộng của mô hình trong quy mô đô thị lớn.

    \item \textbf{Về phạm vi công nghệ và hạ tầng triển khai}\\
    Đề tài sử dụng hệ quản trị cơ sở dữ liệu PostgreSQL kết hợp với phần mở rộng PostGIS để xây dựng kho dữ liệu không gian phục vụ lưu trữ và xử lý dữ liệu giao thông. Quy trình ETL/ELT được phát triển bằng ngôn ngữ Python kết hợp với các thư viện hỗ trợ xử lý dữ liệu và kết nối cơ sở dữ liệu như SQLAlchemy, GeoAlchemy2 và Pandas. Hệ thống được triển khai trên nền tảng điện toán đám mây Microsoft Azure, sử dụng Docker và Docker Compose để quản lý môi trường thực thi và hỗ trợ khả năng mở rộng hệ thống. Các mô hình trí tuệ nhân tạo và học máy được xây dựng và huấn luyện thông qua các thư viện khoa học dữ liệu như Scikit-learn, TensorFlow hoặc PyTorch.

    \item \textbf{Về phạm vi chức năng hệ thống}\\
    Đề tài tập trung nghiên cứu và phát triển nền tảng phần mềm phục vụ thu thập, xử lý, lưu trữ, phân tích và trực quan hóa dữ liệu giao thông. Hệ thống hỗ trợ xây dựng dashboard giám sát giao thông, phân tích dữ liệu lịch sử và tích hợp mô hình dự báo tình trạng ùn tắc giao thông theo thời gian thực. Đề tài không bao gồm việc thiết kế hoặc triển khai hạ tầng phần cứng vật lý như camera giao thông, thiết bị IoT hoặc hệ thống cảm biến ngoài thực địa. Các nguồn dữ liệu đầu vào được giả định là có sẵn thông qua API, dữ liệu mở hoặc các nền tảng hỗ trợ nghiên cứu.

    \item \textbf{Về phạm vi nghiên cứu mô hình dự báo}\\
    Đề tài tập trung vào các bài toán dự báo ngắn hạn liên quan đến lưu lượng giao thông, tốc độ lưu thông và nguy cơ ùn tắc trên từng tuyến đường hoặc khu vực cụ thể. Các mô hình học máy và học sâu được xây dựng nhằm đánh giá khả năng ứng dụng của trí tuệ nhân tạo trong hỗ trợ điều hành giao thông thông minh, tuy nhiên chưa đi sâu vào các bài toán tối ưu điều khiển tín hiệu giao thông hoặc điều phối giao thông quy mô toàn thành phố.

\end{itemize}

\subsection{Ý nghĩa}

\subsubsection{Ý nghĩa thực tiễn}

Đề tài mang ý nghĩa thực tiễn quan trọng trong bối cảnh nhu cầu xây dựng hệ thống giao thông thông minh và đô thị thông minh tại Thành phố Hồ Chí Minh ngày càng trở nên cấp thiết. Việc xây dựng nền tảng kho dữ liệu giao thông tích hợp không chỉ góp phần nâng cao hiệu quả quản lý dữ liệu mà còn tạo tiền đề cho việc ứng dụng các công nghệ phân tích dữ liệu lớn và trí tuệ nhân tạo trong công tác điều hành giao thông đô thị.

\begin{itemize}

    \item \textbf{Hỗ trợ công tác quản lý và điều hành giao thông đô thị}\\
    Hệ thống kho dữ liệu giao thông cho phép tập trung, chuẩn hóa và lưu trữ dữ liệu từ nhiều nguồn khác nhau trong một nền tảng thống nhất, giúp các cơ quan quản lý như Sở Giao thông Vận tải Thành phố Hồ Chí Minh có được cái nhìn toàn diện về hiện trạng giao thông theo không gian và thời gian. Thông qua hệ thống dashboard và các công cụ trực quan hóa dữ liệu, nhà quản lý có thể theo dõi lưu lượng giao thông theo thời gian thực, đánh giá mức độ ùn tắc tại các khu vực trọng điểm và hỗ trợ ra quyết định trong công tác phân luồng, điều phối giao thông cũng như hoạch định chính sách phát triển hạ tầng.

    \item \textbf{Nâng cao hiệu quả khai thác và sử dụng dữ liệu giao thông}\\
    Việc xây dựng quy trình ETL/ELT tự động và kho dữ liệu không gian giúp giải quyết tình trạng dữ liệu giao thông phân tán, thiếu đồng bộ và khó khai thác. Hệ thống cho phép lưu trữ dữ liệu lịch sử với quy mô lớn, hỗ trợ các tác vụ thống kê, phân tích xu hướng và truy vấn dữ liệu phức tạp, từ đó nâng cao hiệu quả sử dụng dữ liệu phục vụ nghiên cứu và vận hành thực tiễn.

    \item \textbf{Hỗ trợ dự báo và giảm thiểu ùn tắc giao thông}\\
    Việc tích hợp các mô hình trí tuệ nhân tạo và học máy vào hệ thống cho phép dự báo lưu lượng phương tiện, tốc độ lưu thông và nguy cơ ùn tắc trong tương lai gần. Các kết quả dự báo có thể hỗ trợ cơ quan quản lý chủ động trong công tác điều tiết giao thông, bố trí lực lượng chức năng hợp lý và đưa ra các giải pháp ứng phó kịp thời nhằm giảm thiểu tác động tiêu cực của ùn tắc giao thông đối với đời sống đô thị.

    \item \textbf{Tạo nền tảng cho phát triển hệ sinh thái giao thông thông minh}\\
    Hệ thống được thiết kế theo hướng kiến trúc mở và có khả năng mở rộng linh hoạt, cho phép tích hợp thêm nhiều nguồn dữ liệu và dịch vụ thông minh khác trong tương lai như dữ liệu IoT, camera giao thông, dữ liệu thời tiết hoặc các nền tảng giao thông công cộng. Điều này góp phần hình thành nền tảng hạ tầng dữ liệu phục vụ phát triển hệ sinh thái giao thông thông minh và đô thị thông minh tại Việt Nam.

    \item \textbf{Khả năng ứng dụng và mở rộng thực tiễn}\\
    Mô hình hệ thống được triển khai trên nền tảng điện toán đám mây và sử dụng các công nghệ mã nguồn mở phổ biến, qua đó bảo đảm tính linh hoạt, khả năng mở rộng và tối ưu chi phí triển khai. Kết quả nghiên cứu có thể được sử dụng làm cơ sở tham khảo cho các bài toán quản lý giao thông tại các đô thị lớn khác ở Việt Nam.

\end{itemize}

\subsubsection{Ý nghĩa khoa học}

Bên cạnh giá trị thực tiễn, đề tài còn mang ý nghĩa khoa học trong lĩnh vực dữ liệu lớn, hệ thống thông tin địa lý và trí tuệ nhân tạo ứng dụng trong giao thông thông minh.

\begin{itemize}

    \item \textbf{Đóng góp về mô hình hóa dữ liệu giao thông không gian – thời gian}\\
    Đề tài nghiên cứu và xây dựng mô hình kho dữ liệu giao thông theo hướng tích hợp dữ liệu không gian – thời gian (spatio-temporal data warehouse), kết hợp giữa dữ liệu giao thông thời gian thực và dữ liệu địa lý không gian. Việc áp dụng các mô hình dữ liệu đa chiều trong bối cảnh dữ liệu giao thông tại Việt Nam góp phần kiểm chứng tính khả thi của phương pháp tiếp cận này đối với các bài toán phân tích giao thông đô thị.

    \item \textbf{Đóng góp trong lĩnh vực xử lý và tích hợp dữ liệu giao thông đa nguồn}\\
    Nghiên cứu đề xuất quy trình ETL/ELT phục vụ thu thập, làm sạch, chuẩn hóa và tích hợp dữ liệu từ nhiều nguồn dị biệt khác nhau. Đây là cơ sở quan trọng cho việc xây dựng các hệ thống dữ liệu giao thông có khả năng xử lý dữ liệu lớn, dữ liệu thời gian thực và dữ liệu không gian trong môi trường đô thị thông minh.

    \item \textbf{Đánh giá khả năng ứng dụng trí tuệ nhân tạo trong dự báo giao thông}\\
    Đề tài thực hiện thử nghiệm các mô hình học máy và học sâu trên dữ liệu giao thông thực tế nhằm đánh giá hiệu quả dự báo lưu lượng và tình trạng ùn tắc giao thông. Kết quả nghiên cứu góp phần bổ sung cơ sở thực nghiệm cho việc ứng dụng trí tuệ nhân tạo trong lĩnh vực giao thông thông minh tại Việt Nam, đặc biệt trong bối cảnh giao thông hỗn hợp và có mật độ phương tiện cao.

    \item \textbf{Đóng góp về kiến trúc hệ thống dữ liệu giao thông thông minh}\\
    Kiến trúc hệ thống được xây dựng theo mô hình phân tầng bao gồm tầng thu thập dữ liệu, tầng kho dữ liệu không gian, tầng trí tuệ nhân tạo và tầng trực quan hóa. Mô hình này có thể được xem như một hướng tiếp cận tham khảo cho việc phát triển các nền tảng Traffic Integrated Operations Center (Traffic IOC) hoặc các hệ thống phân tích giao thông thông minh trong tương lai.

\end{itemize}

\subsection{Đóng góp của đề tài}

Đề tài hướng đến việc xây dựng một nền tảng dữ liệu giao thông thông minh tích hợp nhiều thành phần công nghệ hiện đại, bao gồm kho dữ liệu không gian, xử lý dữ liệu thời gian thực, trí tuệ nhân tạo và hệ thống trực quan hóa phục vụ quản lý giao thông đô thị. Trên cơ sở đó, các đóng góp chính của đề tài được xác định như sau:

\begin{itemize}

    \item \textbf{Đề xuất mô hình kho dữ liệu giao thông tích hợp dữ liệu không gian – thời gian}\\
    Đề tài xây dựng mô hình kho dữ liệu giao thông theo hướng tích hợp dữ liệu đa nguồn kết hợp với dữ liệu địa lý không gian (geospatial data). Mô hình dữ liệu được thiết kế nhằm hỗ trợ lưu trữ, quản lý và phân tích dữ liệu giao thông theo cả hai chiều không gian và thời gian, phục vụ hiệu quả cho các bài toán thống kê, phân tích lưu lượng và dự báo giao thông trong môi trường đô thị thông minh.

    \item \textbf{Xây dựng quy trình ETL/ELT tự động cho dữ liệu giao thông thời gian thực}\\
    Đề tài triển khai quy trình thu thập, xử lý và đồng bộ dữ liệu giao thông theo thời gian thực từ nhiều nguồn dữ liệu khác nhau thông qua API và dữ liệu bản đồ số. Quy trình ETL/ELT được tối ưu nhằm bảo đảm khả năng xử lý dữ liệu liên tục, tự động hóa việc làm sạch dữ liệu, xử lý ngoại lệ và chuẩn hóa dữ liệu trước khi lưu trữ vào kho dữ liệu trung tâm.

    \item \textbf{Ứng dụng công nghệ cơ sở dữ liệu không gian trong quản lý giao thông}\\
    Đề tài khai thác khả năng của PostgreSQL kết hợp với PostGIS để xây dựng kho dữ liệu không gian phục vụ lưu trữ và xử lý các thực thể giao thông như tuyến đường, đoạn đường, nút giao và dữ liệu lưu lượng theo vị trí địa lý. Việc ứng dụng công nghệ dữ liệu không gian góp phần nâng cao khả năng phân tích giao thông theo khu vực và hỗ trợ trực quan hóa dữ liệu trên nền bản đồ số.

    \item \textbf{Tích hợp mô hình trí tuệ nhân tạo phục vụ dự báo giao thông}\\
    Đề tài nghiên cứu và triển khai các mô hình học máy và học sâu nhằm dự báo tình trạng giao thông trong ngắn hạn, bao gồm dự báo lưu lượng phương tiện, tốc độ lưu thông và nguy cơ ùn tắc. Các mô hình được huấn luyện và đánh giá trên tập dữ liệu thực tế đã được xử lý và chuẩn hóa từ kho dữ liệu giao thông, qua đó kiểm chứng khả năng ứng dụng của trí tuệ nhân tạo trong bài toán giao thông đô thị tại Việt Nam.

    \item \textbf{Xây dựng hệ thống trực quan hóa và hỗ trợ ra quyết định}\\
    Đề tài phát triển nền tảng dashboard và web application phục vụ giám sát giao thông theo thời gian thực, hỗ trợ trực quan hóa dữ liệu dưới dạng biểu đồ, bản đồ nhiệt và thống kê lưu lượng giao thông. Hệ thống góp phần hỗ trợ cơ quan quản lý trong công tác điều hành, phân luồng và đánh giá hiện trạng giao thông đô thị.

    \item \textbf{Đề xuất kiến trúc hệ thống giao thông thông minh có khả năng mở rộng}\\
    Đề tài xây dựng kiến trúc hệ thống theo mô hình phân tầng bao gồm tầng thu thập dữ liệu, tầng kho dữ liệu, tầng trí tuệ nhân tạo và tầng giao diện người dùng. Hệ thống được triển khai trên nền tảng điện toán đám mây kết hợp với công nghệ container hóa, qua đó bảo đảm tính linh hoạt, khả năng mở rộng và khả năng tích hợp thêm các nguồn dữ liệu hoặc dịch vụ thông minh trong tương lai.

    \item \textbf{Đóng góp về mặt thực nghiệm trong xử lý dữ liệu giao thông quy mô lớn}\\
    Trong quá trình triển khai, đề tài thực hiện nhiều giải pháp tối ưu hóa hệ thống nhằm đáp ứng yêu cầu xử lý dữ liệu giao thông thời gian thực, bao gồm tối ưu hóa hiệu năng ETL, xử lý đa luồng, quản lý kết nối cơ sở dữ liệu và tối ưu hạ tầng triển khai trên môi trường điện toán đám mây. Các kết quả thực nghiệm góp phần cung cấp cơ sở tham khảo cho việc xây dựng các hệ thống dữ liệu giao thông thông minh trong thực tế.

\end{itemize}
